\section*{Chapitre \ref{ch:m1:basis-roll-and-delivery} --- Base, roll et livraison}

\begin{solution}{exo:m1:basis-roll-and-delivery:1}
$6\,000 \times (1 + 0.042 \times 91/360) = 6\,063.70$; \omterm{def:m1:basis-roll-and-delivery:basis}{base} de 63.70 points d'indice.
\end{solution}

\begin{solution}{exo:m1:basis-roll-and-delivery:2}
$6.0 \times (1 + 0.042 \times 64/360) + 9.5 \times (1 + 0.042 \times 32/360) + 5.5
\times (1 + 0.042 \times 8/360) = 6.04 + 9.54 + 5.51 = 21.09$. \omterm{def:m1:basis-roll-and-delivery:carry}{Fair value}
$6\,063.70 - 21.09 = 6\,042.61$.
\end{solution}

\begin{solution}{exo:m1:basis-roll-and-delivery:3}
\omterm{def:m1:basis-roll-and-delivery:contango}{Contango}~; une position longue roule de 70.00 vers 70.84 chaque mois et, si rien
bouge, voit 70.84 devenir 70.00~: $(70/70.84 - 1) \times 12 = -14.2\%$ par an.
\omterm{def:m1:basis-roll-and-delivery:contango}{Backwardation}~; $(70/69.30 - 1) \times 12 = +12.1\%$ par an.
\end{solution}

\begin{solution}{exo:m1:basis-roll-and-delivery:4}
Le contrat à terme est cher de $6\,050.50 - 6\,042.61 = 7.89$ points d'indice, au-delà de la
borne supérieure de $+6.47$. Vendre le contrat à terme à 6\,050.50~; acheter le panier, 300\,000~\$ par
contrat, dans les proportions de l'indice~; le financer à l'indice de référence majoré de 15
points de base~; encaisser et réinvestir les dividendes~; le matin de l'échéance, vendre chaque
action à l'enchère d'ouverture. Gain verrouillé~: $7.89 - 6.47 = 1.41$ points d'indice, 70~\$ par
contrat. Ce qui peut mal tourner~: le legging (le panier de 500 actions n'est pas acheté
en un seul instant)~; dividendes réduits ou décalés~; hausse du taux de financement~; actions
suspendues à l'ouverture le jour de l'échéance~; \omterm{def:m1:clearing-and-settlement:margin}{marge de variation} sur le contrat à terme vendu
pendant une hausse, à financer alors que le gain du panier n'est pas encore réalisé.
\end{solution}

\begin{solution}{exo:m1:basis-roll-and-delivery:5}
Dichotomie sur $F^*(r) = 6\,047.20$ donne $r = 4.50\%$~: vérification, $6\,000 \times
0.0450 \times 91/360 = 68.25$, dividendes 21.09, $F^* = 6\,047.16$, le reste
dû à l'arrondi du taux. Une hausse de 25 points de base ajoute $6\,000 \times
0.0025 \times 91/360 = 3.79$ points d'intérêt et un montant négligeable de
réinvestissement~: la \omterm{def:m1:basis-roll-and-delivery:carry}{fair value} augmente de 3.79 points, l'indice restant inchangé.
Les contrats à terme sont aussi des instruments de taux.
\end{solution}

\begin{solution}{exo:m1:basis-roll-and-delivery:6}
L'indice publié à 09:30 mélange les actions déjà ouvertes avec la
dernière clôture de celles qui ne le sont pas encore~; il ne vaut 5\,990 que par cette convention.
La cotation utilise le cours d'ouverture propre à chaque action, quand il se produit. Si
les actions qui ouvrent tard ouvrent en baisse, leurs cours d'ouverture entrent dans la
cotation, mais leur contribution au «~plus bas de séance~» de l'indice publié
intervient plus tard, aux cours alors en vigueur des autres actions~: la cotation est un
portefeuille de cours relevés à des instants différents, qui n'ont pas nécessairement à se situer dans la
fourchette d'un indice simultané quelconque. Les détenteurs de contrats à terme et d'options
arrivant à échéance se règlent dessus~; les arbitragistes qui vendent leurs paniers lors des
enchères d'ouverture reçoivent exactement ce cours.
\end{solution}

\begin{solution}{exo:m1:basis-roll-and-delivery:7}
0.751 en \omterm{def:m1:basis-roll-and-delivery:contango}{contango} et 1.273 en \omterm{def:m1:basis-roll-and-delivery:contango}{backwardation} après 24 \omterm{def:m1:basis-roll-and-delivery:roll}{roll} mensuels. Avec un
\omterm{def:m1:basis-roll-and-delivery:contango}{contango} se réduisant linéairement de 1.2~\% à zéro~: 0.861. La moitié de la pente sur les
\omterm{def:m1:pnl-and-positions:realised}{coûts moyens} réduit à peu près de moitié la perte.
\end{solution}

\begin{solution}{exo:m1:basis-roll-and-delivery:8}
La prime est le carry~: trois mois d'intérêts moins trois mois de
dividendes, 42.6 points d'indice à ces taux, et il serait le même si tout le monde
attendait un krach. Les anticipations font bouger ensemble l'indice et le contrat à terme. Ce
que le contrat à terme indique avant l'ouverture, c'est le \emph{changement}~: comparer le
contrat à terme avec la \omterm{def:m1:basis-roll-and-delivery:carry}{fair value} de la veille à la clôture (clôture de la veille plus la
\omterm{def:m1:basis-roll-and-delivery:basis}{base} juste)~; cet écart, et non la prime brute, est le mouvement de nuit
du marché.
\end{solution}

\begin{solution}{pb:m1:basis-roll-and-delivery:1}
\textbf{1.} $1\,000\,000\,000/300\,000 = 3\,333$ contrats.
\textbf{2.} 6\,042.61.
\textbf{3.} 182 jours d'intérêt, 127.40 points d'indice, moins les six dividendes
capitalisés jusqu'au 19 mars, 42.39 points d'indice~: 6\,085.01.
\textbf{4.} 42.39 points d'indice.
\textbf{5.} Intérêts sur les 91 jours entre les échéances sur 6\,000~: 63.70.
Dividendes de ce trimestre, réinvestis~: 21.31. $63.70 - 21.31 = 42.39$.
\textbf{6.} $4.55/6\,000/(91/360) = 30$ points de base par an.
\textbf{7.} $4.55 \times 50 \times 3\,333 = \$758\,258$.
\textbf{8.} $0.025 \times 50 \times 3\,333 = \$4\,166$.
\textbf{9.} Quatre \omterm{def:m1:basis-roll-and-delivery:roll}{roll}~: 3\,049\,695~\$, 30.5 points de base de l'exposition.
\textbf{10.} Les contrats à terme coûtent 30.5 contre 7 pour le fonds, mais la trésorerie du fonds
ne rapporte l'indice de référence dans les deux cas que si le fonds est entièrement financé~; les contrats à terme
laissent 95~\% de la trésorerie libre (pour un fonds qui peut la faire rapporter davantage que l'indice de référence
ou qui en a besoin comme collatéral ailleurs), évitent les écarts de \omterm{def:m1:delta-one-instruments:wht}{retenue à la
source}, et peuvent être vendus à découvert. Pour un détenteur entièrement financé et uniquement acheteur, à ces
niveaux, le fonds indiciel est moins cher.
\textbf{11.} Avec la foule~: c'est alors que le carnet de spreads est le plus profond et le
tick représente un faible coût~; la cherté est ce qu'elle est. Plus tôt~: un détenteur qui
croit que la cherté monte à l'approche de la semaine de \omterm{def:m1:basis-roll-and-delivery:roll}{roll} (à mesure que les bilans des dealers
se remplissent) paie moins en passant le premier, au prix d'un carnet plus mince.
\textbf{12.} Le contrat de mars couvre la fin d'année, lorsque les banques réduisent leurs
bilans aux dates de publication et que le financement est le plus rare~; le \omterm{def:m1:basis-roll-and-delivery:roll}{roll} de décembre
en tient compte.
\textbf{13.} La \omterm{def:m1:basis-roll-and-delivery:basis}{base} elle-même, en points d'indice par rapport à une clôture connue, donc
le \omterm{ex:m1:basis-roll-and-delivery:43}{taux de financement implicite} jusqu'à l'échéance, sans avoir à synchroniser le cours du
contrat à terme avec un indice publié avec retard.
\textbf{14.} Il le peut, car il détient le panier ou peut le financer~: il vend
le contrat à terme cher et touche l'indice de référence majoré de 20. Il le ferait, car 20 points
de base sur une position de bilan relève de son métier~; il ne le fera pas en fin d'année
si le coût du capital dépasse ce niveau.
\textbf{15.} La pente n'est pas un taux d'intérêt mais un rendement de stockage et de
convenance, qui peut atteindre plusieurs pour cent par an, change de signe, et le choix
du contrat à détenir le long de la courbe devient une stratégie en soi.
\textbf{16.} Le E-mini se règle en espèces sur le \omterm{def:m1:basis-roll-and-delivery:soq}{cours spécial
d'ouverture} (SOQ)~: le fonds reçoit ou verse une dernière \omterm{def:m1:clearing-and-settlement:margin}{marge de variation} et est hors
du marché, sans couverture, dès ce matin-là. Le WTI est livré~: un acheteur qui n'a pas
clôturé est assigné en barils à Cushing et doit avoir prévu un transport par pipeline ou
un stockage~; les courtiers clôturent d'office ces positions avant d'en arriver là.
\textbf{17.} C'est le prix de marché du bilan des dealers pour le financement en
actions, le même chiffre que celui qui figure dans les spreads de swaps, dans le coût des
stratégies à effet de levier et dans les rendements de l'arbitrage sur indices.
\textbf{18.} Celui qui est vendeur du contrat à terme et acheteur de l'action~: les
arbitragistes d'indice, dealers et investisseurs financés qui remplacent leurs actions par des
contrats à terme lorsque le \omterm{def:m1:basis-roll-and-delivery:roll}{roll} est bon marché, et l'inverse lorsqu'il est cher.
\textbf{19.} \textbf{30.5 points de base par an.}
\textbf{20.} Parce qu'à l'échéance il verse l'indice, donc quiconque peut détenir le
panier jusqu'à cette date peut encaisser tout écart.
\end{solution}

\begin{solution}{iq:m1:basis-roll-and-delivery:1}
L'acheteur du contrat à terme ne paie pas les actions et perçoit des intérêts sur
la trésorerie, mais ne reçoit aucun dividende~; la \omterm{def:m1:basis-roll-and-delivery:carry}{fair value} vaut le cours au comptant
plus les intérêts moins les dividendes. Il se traite avec une décote par rapport au
comptant lorsque le rendement en dividendes jusqu'à l'échéance dépasse le taux d'intérêt,
comme les années de taux bas ou sur les marchés où la saison des dividendes est dense avant
l'échéance, et, par rapport à la \omterm{def:m1:basis-roll-and-delivery:carry}{fair value}, lorsque le panier est coûteux à vendre à découvert.
\iqlookfor{le carry, non les anticipations~; un cas concret de décote.}
\end{solution}

\begin{solution}{iq:m1:basis-roll-and-delivery:2}
\omterm{def:m1:basis-roll-and-delivery:contango}{Contango}~: les échéances éloignées au-dessus des plus proches~; \omterm{def:m1:basis-roll-and-delivery:contango}{backwardation}~: en dessous.
Aucune des deux n'est une prévision. Pour un actif financier, la pente vaut taux moins rendement~;
pour une matière première, elle ajoute le coût de stockage et retranche le rendement de commodité.
Une courbe en \omterm{def:m1:basis-roll-and-delivery:contango}{contango} signifie que porter la matière première coûte cher et que les stocks sont abondants~;
l'investisseur acheteur supporte la pente sous forme de rendement de \omterm{def:m1:basis-roll-and-delivery:roll}{roll} négatif.
\iqlookfor{la pente comme carry, et la conséquence sur le rendement de \omterm{def:m1:basis-roll-and-delivery:roll}{roll}.}
\end{solution}

\begin{solution}{iq:m1:basis-roll-and-delivery:3}
Calculer la \omterm{def:m1:basis-roll-and-delivery:carry}{fair value} à partir du taux et des dividendes datés. Lorsque le contrat à terme
dépasse la \omterm{def:m1:basis-roll-and-delivery:carry}{fair value} de plus que mes coûts, vendre des contrats à terme et acheter le panier
par un programme de transactions, financé. Carry~: encaisser les dividendes, payer le
financement, payer ou recevoir la \omterm{def:m1:clearing-and-settlement:margin}{marge de variation}. Sortir soit lorsque l'écart se résorbe
(vendre le panier, acheter le contrat à terme), soit à l'échéance, en vendant chaque action
dans sa fixing d'ouverture le matin de l'échéance, ce qui réalise exactement le \omterm{def:m1:basis-roll-and-delivery:soq}{cours spécial
d'ouverture}. Risques~: legging, dividendes, financement, suspensions de cotation, trésorerie de marge.
\iqlookfor{la sortie par la fixing d'ouverture et la trésorerie de marge sur le contrat à terme vendu.}
\end{solution}

\begin{solution}{iq:m1:basis-roll-and-delivery:4}
(i) Échéance la plus proche non ajustée~: véritables cours traités, avec un saut à chaque \omterm{def:m1:basis-roll-and-delivery:roll}{roll}~; à utiliser
pour les niveaux à une date donnée, jamais pour les rendements. (ii) Ajustement par différence
(rétro-ajustement par l'écart à chaque \omterm{def:m1:basis-roll-and-delivery:roll}{roll})~: conserve le P\&L en points, déforme les rendements
en pourcentage et peut devenir négatif. (iii) Ajustement par ratio~: conserve les rendements
en pourcentage, déforme les niveaux. (iv) Le meilleur choix~: calculer les rendements à partir
des cours du même contrat et les enchaîner, en conservant l'identité du contrat~; dériver ensuite
toute série de niveaux à partir de celle-ci. La règle de date de \omterm{def:m1:basis-roll-and-delivery:roll}{roll} (croisement des volumes ou
des \omterm{def:m1:futures-contracts-and-exchanges:oi}{positions ouvertes}, ou nombre de jours fixe) doit être explicite et doit pouvoir être suivie
par la stratégie.
\iqlookfor{rendements à partir des cours du même contrat~; conscience que les ajustements par ratio
et par différence répondent à des questions différentes.}
\end{solution}

\begin{solution}{iq:m1:basis-roll-and-delivery:5}
Le contrat est \omterm{def:m1:basis-roll-and-delivery:delivery}{à livraison physique} sur un hub sans accès à la mer. La veille de l'échéance, les détenteurs
de positions longues qui ne pouvaient pas prendre livraison ont dû vendre, le stockage du hub était engagé,
et ceux qui pouvaient recevoir du pétrole exigeaient une rémunération pour le faire~; comme une grande partie
du volume de la journée était liée au règlement-livraison par des ordres \omterm{def:m1:basis-roll-and-delivery:tas}{trade at settlement}, la fenêtre de
règlement-livraison présentait des vendeurs et presque aucun acheteur. Un contrat à terme sur indice \omterm{def:m1:basis-roll-and-delivery:delivery}{réglé en espèces}
ne peut pas produire cet effet~: son prix final est l'indice, qui est positif ou nul, et personne n'est jamais obligé
de stocker quoi que ce soit.
\iqlookfor{livraison et stockage comme mécanisme~; le règlement en espèces
comme différence.}
\end{solution}

\begin{solution}{iq:m1:basis-roll-and-delivery:6}
Vendre le \omterm{def:m1:basis-roll-and-delivery:roll}{roll}~: acheter décembre, vendre mars, contre un panier acheteur ou, plus simplement, être vendeur
de contrats à terme de mars contre des actions détenues à partir de l'échéance de décembre, en
gagnant l'indice de référence plus 60 points de base sur la période couvrant la fin d'année.
Formes équivalentes~: vendre des BTIC, prêter par un \omterm{def:m1:financing:trs}{total return swap}. Limites~: la position consomme
du bilan et du capital à cheval sur la date de publication, ce qui explique précisément pourquoi elle est
chère~; la position est longue de \$X d'actions financées sur le passage de fin d'année~; \omterm{def:m1:pnl-and-positions:mtm}{valorisation au prix
de marché} du spread si la cherté s'accentue~; risque de dividende~; et, pour un fonds, la trésorerie de marge
sur le contrat à terme vendu. Dimensionner la position selon le bilan que l'on peut engager à cheval sur
le passage de fin d'année, et non selon le spread.
\iqlookfor{reconnaître que la cherté est le prix de la contrainte
qui limite la transaction.}
\end{solution}
